\documentclass[11pt,a4paper]{article}
\usepackage[T1]{fontenc}
\usepackage[utf8]{inputenc}
\usepackage[english]{babel}
\usepackage{lmodern}
\usepackage[margin=27mm]{geometry}
\usepackage{microtype}
\usepackage[hidelinks]{hyperref}
\title{Cross-Sensor and Cross-Domain Transfer of RGB and Long-Wave Infrared Detectors}
\author{Mateusz Pomianek}
\date{Version 0.1.0 -- 21 November 2026}
\begin{document}
\maketitle

\begin{abstract}
Detection performance may degrade across genuinely independent sensor platforms, viewpoints and thermal domains. Can training-only bounded adaptation reduce detection error on unseen RGB/LWIR sensors without access to held-out target annotations? The prospective design is as follows: Choose independent legally reusable sensor-domain cohorts with harmonized class definitions and explicit optics, resolution and spectral-band metadata; freeze adaptation on training cohorts. The primary outcome is Out-of-domain recall loss at a matched false-positive operating point on independently held-out sensor cohorts. The manuscript defines a research question and an evaluation strategy; it does not report experimental estimates.
\end{abstract}

\section{Introduction}
Detection performance may degrade across genuinely independent sensor platforms, viewpoints and thermal domains. The scientific contribution under consideration is a falsifiable comparison in which alternative explanations, measurement limitations and distinct hypotheses from related ARC Robo studies can be assessed independently. Both the source of observations and the assumptions needed to interpret results must be documented before any confirmatory comparison.

\section{Research question and hypothesis}
Can training-only bounded adaptation reduce detection error on unseen RGB/LWIR sensors without access to held-out target annotations? The working hypothesis is that the proposed approach improves the primary endpoint under an explicitly matched evidence and resource budget. This hypothesis is subject to prior-art review and prospective specification; non-superiority and negative findings must remain reportable.

\section{Materials and methods}
Choose independent legally reusable sensor-domain cohorts with harmonized class definitions and explicit optics, resolution and spectral-band metadata; freeze adaptation on training cohorts. The experimental unit, independent ground truth, data inclusion criteria, version of software and sensors, and provenance of each dataset or scenario will be fixed in the protocol. Training, calibration and evaluation partitions will be isolated at the appropriate physical scene, sequence, participant or device level.

\subsection{Comparators and experimental controls}
Compare source-only inference, controlled fine-tuning and explicitly parameter-limited adaptation on unseen target-device sequences. Model or policy hyperparameters must not be selected from final held-out results. The implementation will record relevant software versions, seeds, sensor configuration, timestamps and data rights.

\subsection{Primary endpoint}
Out-of-domain recall loss at a matched false-positive operating point on independently held-out sensor cohorts. Denominators, operational thresholds, missing observations and failure cases will be defined before inferential analysis. Secondary metrics will not be substituted post hoc for the primary endpoint.

\section{Statistical and computational analysis plan}
Estimate domain-wise generalization gaps and uncertainty over source/target sensor pairs, keeping adaptation and evaluation data disjoint. Any confirmatory statistical contrast will report uncertainty intervals, effective independent sample size and relevant failure strata. Where a physical validation or participant-based experiment is necessary, approval and feasibility assessment must precede data collection.

\section{Anticipated interpretation and limitations}
Image resizing of one sensor does not constitute cross-sensor transfer; incompatible labels or infrared bands cannot be silently pooled. Confounding from software versions, sensor conditions, scene selection and methodological overlap will be assessed explicitly. Claims about simulated behavior will not be generalized to physical safety or human behavior without matching evidence.

\section{Results}
Empirical results, tables and uncertainty intervals will be incorporated after execution of the documented evaluation protocol and independent verification of the input evidence.

\section{Discussion}
The planned comparison is intended to determine whether the proposed mechanism yields an identifiable gain over the stated controls. Interpretation must distinguish failure of a scientific hypothesis from an incomplete implementation, and should describe practical operating conditions under which conclusions remain valid.

\section{Conclusion}
This study defines a testable research problem, matched comparators and a primary evaluation target. A substantive empirical conclusion remains conditional on reproducible measurements and appropriate methodological review.

\section{Data and code availability}
The authoritative scientific source is maintained in the \texttt{MatPomGit/arc-rgbt} repository. Dataset permissions, input manifests, source commits, evaluation scripts and permitted public artifacts will be identified in the final version. No participant records or proprietary instrumentation details are included here.
\end{document}
